\section{Scope and Limitations}\label{sec:limits}

We adopt the same posture as the FTRFS v1 report: the limitations
section is structural, not concessionary. It records, in plain
language, what the present version covers and what it does not.

\paragraph{V1 covers.}
The mathematical formalization (state space, invariants, recovery
operator, soundness theorem with proof in \cref{app:proof}), the
algorithmic derivation, the implementation-path mapping, and a
re-reading of the FTRFS v1 dataset under the calculus. Three
memory-stack layers are modelled physically: DRAM, SRAM, NAND flash
via NVMe. Sources are canonical: JEDEC, NASA-HDBK, ESA ECSS,
IEEE TNS.

\paragraph{V1 does not cover.}
\begin{itemize}
\item \emph{Beam-time measurements.} No accelerator data; the surface
  of \cref{fig:recovery-surface} is a projection from canonical
  cross-section data, explicitly so.
\item \emph{Verified-compilation closure.} The implementation path
  preserves invariants by construction, but the gap from
  algorithm to kernel binary is the standard
  unverified-compilation gap. seL4-style closure is v3 work.
\item \emph{TID and DDD modelling.} Cumulative ionizing dose and
  displacement damage are below the filesystem layer; they are
  hardware-level concerns. Mention only.
\item \emph{Adversarial fault injection.} An attacker with kernel
  privileges who manipulates the filesystem state directly is
  outside~\(\mathcal{E}\). This is the same scope choice as
  FTRFS v1 (Stage~5 in that roadmap).
\item \emph{Hardware SEL recovery.} Latched-up hardware is recovered
  by power cycling; the filesystem operator's response is to map
  it to~\(\bot\), nothing more.
\item \emph{Emerging memory technologies.} 3D~XPoint, ReRAM, MRAM
  are positioned but not modelled.
\item \emph{DDR5 on-die ECC.} Mentioned as a strengthening
  assumption to be added in v2; v1 models the unfiltered DRAM
  matrix as the conservative baseline.
\item \emph{The fault-injection metric M6.} Reserved for v2,
  following the FTRFS Phase~E campaign currently under
  development.
\end{itemize}

\paragraph{Methodological limits.}
The author is the sole contributor; no peer review beyond the
public mailing-list interactions on linux-fsdevel
(referenced in~\cite{ftrfs-rfc-thread}) has shaped this text.
The proof of the soundness theorem (\cref{app:proof}) is
hand-written; it is not machine-checked. v3 will explore
mechanization in a proof assistant such as Coq, Isabelle/HOL,
or Lean.

\paragraph{Threats to validity.}
The principal threat is that the perturbation family~\(\mathcal{E}\)
defined in \cref{sec:phys} does not faithfully represent the
distribution of real SEE events seen by a deployed filesystem.
This is precisely why beam-time measurements are needed for v2:
to bind the model to ground truth. The soundness theorem is
robust under tightening of~\(\mathcal{E}\) (any sub-family
of~\(\mathcal{E}\) inherits the result) but fragile under
loosening (a perturbation that exceeds RS capacity is, by
construction, out of scope).
